\bookchapter{Kapitel 22: Wie KI im Handy arbeitet}

\chapterintro{Willkommen auf der coolsten „Magie-Bühne“ des Smartphones – der Hauptdarsteller ist: Künstliche Intelligenz. Die Sofort-Fotos, die du täglich machst, der Sprachassistent, der immer parat ist, die Eingabemethode, die dich mit der Zeit immer besser versteht – all das wird von diesem unsichtbaren Gehirn gesteuert.}

\booksection{I. KI-Kamera: Deine „digitale Dunkelkammer“ im Smartphone}

Die heutige KI-Fotografie ist eher so, als würde ein professioneller Retuscheur plus Fotografen-Team direkt in den Auslöser gesteckt. Bevor du abdrückst, arbeitet die KI bereits: Sie scannt die Szene und erkennt: „Das ist blauer Himmel“, „Das ist ein Gesicht”, „Das ist eine Nachtszene“.

\booksection{II. Sprachassistent: Wie dein „Ton-Sekretär“ dich versteht}

\begin{enumerate}
  \item \textbf{Hören (Spracherkennung)}: Schallwellen in Text umwandeln.
  \item \textbf{Verstehen (natürliches Sprachverstehen)}: Intention verstehen, wichtige Entitäten extrahieren.
  \item \textbf{Ausführen und antworten}: Entsprechende App oder Dienst aufrufen und per Sprachsynthese „zu dir sprechen".
\end{enumerate}

\booksection{III. On-Device-KI vs. Cloud-KI}

On-Device-KI behält die Berechnung im Smartphone – schnell, privatsphärenfreundlich, offline-fähig. Cloud-KI benötigt Netz, kann aber komplexere Aufgaben bewältigen.

\booksection{IV. Große KI-Modelle und das Smartphone: Dein intelligenter Begleiter}

2025-2026 hat der Aufstieg grosser KI-Modelle (wie DeepSeek R1, ChatGPT-5, ERNIE 5.0, Qwen 3 usw.) die Art, wie wir Smartphones nutzen, tiefgreifend veraendert. Das Handy ist nicht mehr bloss ein ausfuehrendes Werkzeug, sondern ein intelligenter Begleiter, der verstehen, schliessen und kreieren kann. On-Device-KI-Chips (NPU-Rechenleistung ueber 80 TOPS) ermoeglichen erstmals das lokale Ausfuehren grosser Modelle auf dem Handy – KI funktioniert auch ohne Internetverbindung.

\booksubsection{Was ist ein großes KI-Modell?}

Ein großes KI-Modell (Large Language Model, LLM) ist ein KI-System, das mit riesigen Textmengen trainiert wurde und natürliche Sprache verstehen und erzeugen kann. Stell dir einen „Super-Streber“ vor, der fast alle Texte des Internets gelesen hat – er kann Fragen beantworten, Artikel schreiben, übersetzen, programmieren, Daten analysieren.

\booksubsection{Anwendungen großer KI-Modelle auf dem Smartphone:}

\begin{itemize}
  \item \textbf{Dialog-Assistent}: Apps wie DeepSeek, ChatGPT, Doubao, Kimi – jederzeit Fragen stellen, übersetzen, Texte schreiben, Artikel zusammenfassen. Intelligenter als eine klassische Suchmaschine – sie liefert nicht nur Links, sondern direkte Antworten.
  \item \textbf{Schreiben und Kreatives}: KI hilft beim Schreiben von E-Mails, Zusammenfassen von Besprechungsprotokollen, Überarbeiten von Texten, Generieren von Werbe-Texten.
  \item \textbf{Bilderzeugen und -bearbeitung}: KI entfernt mit einem Tipp Passanten aus Fotos, erweitert Bildränder, wandelt Stile um und erzeugt sogar Bilder aus Textbeschreibungen.
  \item \textbf{Echtzeit-Übersetzung}: Telefonate in Echtzeit übersetzen, Face-to-Face-Übersetzung, Foto-Übersetzung – auf Reisen keine Sprachbarrieren mehr.
  \item \textbf{Lern-Nachhilfe}: KI kann Konzepte erklären, Übungsaufgaben erstellen, Hausaufgaben korrigieren – ein Nachhilfelehrer für die Hosentasche.
\end{itemize}

\booksubsection{On-Device-Großmodelle: Das Smartphone führt KI selbst aus}

Die neuesten Flaggschiffe können kleine KI-Großmodelle (rund 7B Parameter) bereits lokal ausführen. Das bringt:

\begin{itemize}
  \item \textbf{Kein Netz nötig}: KI-Funktionen auch ohne Netz – im Flugzeug, in der U-Bahn.
  \item \textbf{Mehr Privatsphäre}: Deine Dialoge und Daten werden nicht in die Cloud hochgeladen, sondern vollständig lokal auf dem Smartphone verarbeitet.
  \item \textbf{Schnellere Antwort}: Keine Netzwerklatenz – die KI-Antwort erfolgt nahezu sofort.
  \item \textbf{Aktuelle Grenzen}: On-Device-Modelle sind noch nicht so leistungsfähig wie Cloud-Großmodelle – komplexe Aufgaben benötigen weiterhin Netz. Mit steigender KI-Rechenleistung der Chips werden On-Device-Modelle aber immer stärker.
\end{itemize}

\booksubsection{Systemweite KI-Integration}

\begin{itemize}
  \item \textbf{Apple Apple Intelligence}: Das in iOS 18 eingeführte KI-System – integriert Schreibwerkzeuge, Bild-Zauberstab, Siri-Erweiterung, Benachrichtigungs-Zusammenfassung u. a. On-Device-Verarbeitung hat Vorrang, komplexe Aufgaben gehen in die Cloud.
  \item \textbf{Huawei Pangu-Großmodell}: In HarmonyOS integriert – der XiaoYi-Assistent erhält Großmodell-Fähigkeiten, unterstützt Dokument-Zusammenfassung, intelligente Bildschirkerkennung, KI-Bildbearbeitung u. a.
  \item \textbf{Xiaomi Pengpai-KI}: Systemweite KI-Fähigkeiten – unterstützt KI-Bildsuche, KI-Zusammenfassung, KI-Echtzeit-Untertitel, On-Device-Großmodell-Dialog u. a.
  \item \textbf{Samsung Galaxy AI}: Echtzeit-Telefonübersetzung, KI-Notiz-Zusammenfassung, Bild-Erzeugung/-Bearbeitung, Circle-to-Search u. a.
\end{itemize}

\booksection{V. Fünf typische Anwendungsszenarien mit Praxis}

Damit du die „Superkräfte“ der Smartphone-KI anschaulicher erlebst, zeigen wir in fünf konkreten Szenarien, wie die KI im Alltag „zur Rettung eilt”:

\booksubsection{Szenario 1: Echtzeit-Übersetzung – der persönliche Dolmetscher auf Reisen}

\textbf{Tatsächliche Leistung}: In Paris die Speisekarte auf Französisch nicht verstanden? Die iPhone-App „Übersetzung“ unterstützt einen Face-to-Face-Zweisprachen-Modus – beide sehen die Untertitel in ihrer eigenen Sprache, sogar mit Simultanübersetzung. Auf der Android-Seite bietet Samsung Galaxy AI eine Echtzeit-Anrufübersetzung – du sprichst Chinesisch ins Telefon, der Gesprächspartner hört Englisch, und umgekehrt.

\textbf{Dahinterliegendes Prinzip}: Spracherkennung (ASR) maschinelle Übersetzung (NMT) Sprachsynthese (TTS) als Drei-Schritt-Pipeline – die End-to-End-Latenz lässt sich inzwischen auf unter 1 Sekunde kontrollieren.

\booksubsection{Szenario 2: KI-Entfernung – gerettete Urlaubsfotos, die von Passanten „ruiniert“ wurden}

\textbf{Tatsächliche Leistung}: Am Sehenswürdigkeit ein fantastisches Selfie geschossen, aber der Hintergrund ist voller anderer Touristen? Mit „Magie-Entfernung“ die Passanten umranden – das Smartphone erkennt sie automatisch, „malt” sie weg und füllt den Hintergrund intelligent auf. Das Ergebnis ist oft verblüffend gut.

\textbf{Dahinterliegendes Prinzip}: Ein Bild-Inpainting-Modell erkennt und löscht zuerst den markierten Bereich und „denkt“ sich dann basierend auf umliegenden Pixeln und semantischem Verständnis plausible Füllinhalte aus. Das zugrundeliegende KI-Modell wurde mit Millionen von Bildern trainiert.

\booksubsection{Szenario 3: KI-Bildbearbeitung – mit einem Klick zum „Fotografie-Meister“}

\textbf{Tatsächliche Leistung}: Ein trüb-graues Landschaftsfoto an einem bewölkten Tag? KI-„Magischer Himmel“ tauscht mit einem Tipp den Himmel gegen blauen Himmel, Abendrot oder sogar Sternenhimmel aus. KI-Bild-Erweiterung macht aus einem Hochformat ein Querformat – das Smartphone „denkt” sich die Szene außerhalb des Bildes aus.

\textbf{Dahinterliegendes Prinzip}: Auf generativer KI (Generative AI) basierende Diffusionsmodelle (Diffusion Model) verstehen den semantischen Inhalt des Bildes und erzeugen neue Pixel – das ist die „vorderste Front“ der aktuellen KI-Bildverarbeitung.

\booksubsection{Szenario 4: Dokumentverarbeitung – der Effizienz-Booster für Büroarbeiter}

\textbf{Tatsächliche Leistung}: Nach dem Meeting das Audiomaterial der KI übergeben – sie erstellt automatisch strukturierte Besprechungsprotokolle, extrahiert To-dos und generiert eine Zusammenfassung. Ein 10-seitiger englischer PDF-Bericht kann von der KI in 30 Sekunden auf die Kernergebnisse destilliert werden.

\textbf{Dahinterliegendes Prinzip}: Ein Long-Context-Großmodell (Long Context LLM) kann Dokumente mit Zehntausenden von Wörtern „komplett lesen“, mehrstufig schlussfolgern und Informationen extrahieren. Apple Intelligence, Huawei XiaoYi, WPS AI, Tencent Yuanbao u. a. verfügen über diese Fähigkeit.

\booksubsection{Szenario 5: KI-Bildsuche und -Erkennung – deine „visuelle Suchmaschine“}

\textbf{Tatsächliche Leistung}: Auf der Straße eine unbekannte Blume gesehen – Kamera an, und die KI sagt sofort: „Das ist eine Hibiskus.“ In der Galerie das Foto vom letzten Jahr am Meer gesucht – Suche nach „Meer”, und die KI findet automatisch alle Fotos mit Meer, ohne dass du manuell blättern musst.

\textbf{Dahinterliegendes Prinzip}: Ein Computer-Vision-Modell (CV) wandelt den Bildinhalt in Vektoren (Embedding) um und gleicht sie mit der Textanfrage auf Ähnlichkeit ab. Diese „multimodale Suche“ ist ein großer Sprung in der Smartphone-Galerie-Verwaltung.

\booksection{VI. Prompt-Optimierung: Lass die KI nach deiner Pfeife tanzen – das Kommunikations-Prinzip}

Die Qualität der KI-Antworten hängt stark davon ab, wie du „fragst“. Dieselbe Frage bringt bei unterschiedlicher Formulierung völlig unterschiedliche Antworten. Im Folgenden stellen wir dir den\textbf{praktischsten und am leichtesten anzuwendenden „Rolle-Hintergrund-Aufgabe-Grenzen“-Kommunikationsansatz}, begleitet von zahlreichen Beispielen.

\begin{tcbinfo}{Warum ein Framework?}
Viele nutzen KI wie Baidu – sie tippen ein paar Stichworte ein. Aber KI ist keine Suchmaschine, sondern ein Sprachmodell, das nach dem Lesen und Verstehen antwortet.\textbf{Gib ihr mehr Kontext, sage klar, was du willst, und halte sie davon ab, vom Thema abzudriften}– dann liefert sie präzise Ergebnisse. Das Framework übersetzt das „Wie soll die gute Antwort aussehen“ in deinem Kopf in Anweisungen, die die KI versteht.

\end{tcbinfo}

\booksubsection{Rolle (Role): Sag der KI, wer sie ist}

Gib der KI eine konkrete „Identität“, und sie ruft das Fachwissen und die Ausdrucksweise des entsprechenden Bereichs auf.\textbf{Je spezifischer die Rolle, desto professioneller die Ausgabe}.

\booksubsection{Praktische Beispiele:}

\begin{itemize}
  \item „Du bist ein\textbf{Smartphone-Reparaturtechniker mit 10 Jahren Xiaomi-Erfahrung}.“ (technische Frage)
  \item „Du bist ein\textbf{Physik-Lehrer am Gymnasium, der komplexe Konzepte einfach erklärt}." (Wissenschafts-Popularisierung)
  \item „Du bist ein\textbf{Buch-Lektor mit gutem Geschmack, vielseitiger Lektüre und eigenem Stil}." (Schreibproblem)
  \item „Du bist ein\textbf{Koch-Blogger, der jeden Morgen für Berufstätige frühstückt}." (Lebensproblem)
  \item „Du bist ein\textbf{geduldiger, sanfter Arzt, der nie über Anfänger lacht}.“ (Gesundheitsfrage)
\end{itemize}

\textbf{Warum wirkt das?}Die Rollenvorgabe bestimmt\textbf{Wissensbereich},\textbf{Ausdrucksstil}und\textbf{Fachtiefe}. Eine KI ohne Rolle ist ein Generalist, der von allem etwas und nichts richtig kann; mit einer Rolle kann sie in diese Identität „eintauchen“ und wie diese Fachperson denken.

\booksubsection{Hintergrund (Background): Sag der KI deine Situation}

Wer du bist, dein aktueller Stand, deine Einschränkungen – das vorab klar sagen, dann kann die KI\textbf{wirklich zu dir passende}Antworten geben, statt allgemein daherzureden.

\booksubsection{Praktische Beispiele:}

\begin{itemize}
  \item „Mein Smartphone ist ein\textbf{Xiaomi 14, seit einem halben Jahr, wird beim Spielen in letzter Zeit sehr heiß}." (Gerätesituation)
  \item „Ich bin\textbf{eine 60-jährige Rentnerin, habe schlechte Augen und kann nicht gut mit dem Smartphone umgehen}." (Nutzerprofil)
  \item „Ich\textbf{bereite mich auf die Englisch-CET-6-Prüfung im Juni 2026 vor, meine Grundlage ist schwach}.“ (Lernhintergrund)
  \item „Ich habe\textbf{ein kleines Café in einer Drittklass-Stadt, monatlicher Umsatz 30.000 Yuan, zwei Angestellte}." (Geschäftshintergrund)
  \item „Ich\textbf{habe gerade ein Baby bekommen und bin im Wochenbett, mein Schlaf ist stark gestört}." (Lebensphase)
\end{itemize}

\textbf{Warum wirkt das?}Hintergrundinformationen sagen der KI,\textbf{„für wen“}und\textbf{„unter welchen Bedingungen”}. Dasselbe Problem wird Schülern und Eltern unterschiedlich erklärt; Stadtbewohnern und Dorfbewohnern werden unterschiedliche Empfehlungen gegeben. Fehlt der Hintergrund, kann die KI nur „kleinsten gemeinsamen Nenner“-Antworten geben – oft nichtssagend.

\booksubsection{Aufgabe (Task): Sag der KI klar, was sie tun soll}

In diesem Schritt\textbf{Verben aktivieren}(schreiben, auflisten, vergleichen, übersetzen, erklären, zusammenfassen …) und die Aufgabe\textbf{konkret und ausführbar}beschreiben. Vage Aufgaben führen nur zu vagen Antworten.

\booksubsection{Praktische Beispiele:}

\begin{itemize}
  \item „\textbf{Liste 5 Berechtigungen auf, die ich sofort deaktivieren sollte}, und erkläre warum." (Klare Aktion)
  \item „\textbf{Übersetze die folgende englische E-Mail ins Chinesische}, im geschäftlichen Ton." (Konkrete Aufgabe)
  \item „\textbf{Vergleiche diese vier Smartphones um 3000 Yuan}und gib eine Vor-/Nachteile-Tabelle." (Klares Format)
  \item „\textbf{Erkläre „Was ist 5G" so, dass es ein Grundschüler versteht}." (Klarer Stil)
  \item „\textbf{Schreibe drei Kurzvideo-Texte für mein Frühstücksgeschäft}– jeweils für Berufstätige, Paare und Eltern mit Kindern." (Klare Segmentierung)
\end{itemize}

\textbf{Warum wirkt das?}Der Aufgaben-Teil ist das „Befehlszentrum“ für das Handeln der KI.\textbf{Je klarer die Verben und je konkreter die Anforderungen, desto besser weiß die KI, was sie liefern soll}Vage „mach mal“, „erstell mal”, „schau mal“ sind KI-Tabu – sie rät herum und das Ergebnis driftet oft ab.

\booksubsection{Einschränkungen (Constraints): Sag der KI, was sie nicht tun soll}

„Einschränkungen“ klingen wie „Zwänge“, sind in Wahrheit\textbf{das Werkzeug, mit dem die KI auf die für dich passende Ausgabe fokussiert wird}. Häufige Dimensionen: Wortzahl, Format, Tonfall, Vermeidungen, Pflicht-Inhalte.

\booksubsection{Praktische Beispiele:}

\begin{itemize}
  \item „\textbf{Maximal 200 Wörter, als Aufzählung}." (Länge/Format)
  \item „\textbf{Freundlicher, natürlicher Ton wie bei einem Gespräch unter Freunden}, keine Amtssprache wie „erstens, zweitens, drittens“.“ (Tonfall)
  \item „\textbf{Keine Fachbegriffe verwenden}, bei Fachwörtern bitte gleichzeitig erklären." (Wort-Einschränkung)
  \item „\textbf{Maximal 3 Sätze pro Absatz}, für mobiles Lesen." (Layout)
  \item „\textbf{Muss enthalten: Nutzungsszenario, Vor-/Nachteile, Passung zu meinem Budget}.“ (Pflicht-Inhalte)
  \item „\textbf{Keine konkreten Anlageempfehlungen}, nur Prinzipien erklären." (Grenze)
\end{itemize}

\textbf{Warum wirkt das?}Ohne Einschränkungen neigt die KI zum „Freistil“ – zu akademisch, zu lang oder am Thema vorbei. Klare Einschränkungen = klares „Lieferstandard”, beide Seiten haben gleiche Erwartungen, das Ergebnis ist kontrollierbar.

\booksubsection{Prompt-Qualitätsprüfung: 5 Merkmale eines guten Prompts}

Einen guten Prompt zu schreiben, ist keine einmalige Sache –\textbf{nach dem Schreiben diese fünf Punkte selbst prüfen}, bevor du sie an die KI schickst:

\begin{enumerate}
  \item \textbf{Konkret statt abstrakt}: Nicht „schreib einen Artikel", sondern „schreib einen 800-Wörter-Artikel mit 3 Abschnitten, ein Einkaufsleitfaden für Smartphones, für frischgebackene Mütter".
  \item \textbf{Mit Hintergrund statt isoliert}: Sag der KI „wer ich bin, warum ich das frage" – nur dann kann sie „maßgeschneiderte" Antworten geben.
  \item \textbf{Klares Ausgabeformat}: Tabelle, Liste, Absatz oder JSON? Vorab sagen.
  \item \textbf{Einschränkungen präzise}: Wortzahl, Stil, Vermeidungen – drei davon bringen schon deutliche Verbesserung.
  \item \textbf{Von der KI überprüfbar}: Nach dem Fragen die KI bitten, deine Anforderungen einmal zu wiederholen – wenn sie das nicht klar sagen kann, war es nicht klar gesagt.
\end{enumerate}

\booksubsection{Häufige Probleme und Lösungen}

\begin{center}
\begin{tabularx}{\linewidth}{|X|X|}
\hline
\textbf{Symptom} & \textbf{Lösung} \tabularnewline
\hline
Antwort zu allgemein, unbrauchbar & Hintergrund ergänzen (wer du bist, in welchem Szenario) \tabularnewline
Antwort zu akademisch / zu bürokratisch & Tonfall-Einschränkung ergänzen + Liste abgelehnter „Phrasen" \tabularnewline
Antwort am Thema vorbei / Frage nicht beantwortet & Aufgabe mit Verben neu formulieren und verlangen: „nur das beantworten" \tabularnewline
Antwort zu lang, nicht zu Ende lesbar & Hart Wortzahl begrenzen (z. B. unter 200 Wörter) \tabularnewline
KI redet ernsthaft daher & „Wenn unsicher, markiere als ‚ich bin unsicher'" + wichtige Fakten selbst prüfen \tabularnewline
Erste Version nicht zufriedenstellend & Nicht neu schreiben – an den Original-Prompt anhängen: „nochmal überarbeiten" \tabularnewline
\hline
\end{tabularx}
\end{center}

\booksubsection{Drei vollständige Szenario-Beispiele}

\begin{tcbinfo}{Szenario 1: Lebens-Assistent (Frühstücks-Kombinations-Vorschläge)}
\textbf{Prompt}:

Du bist ein Frühstücks-Blogger, der sich auf nährstoffreiche Kombinationen versteht (\textbf{Rolle}). Ich bin eine berufstätige Mutter, die jeden Morgen nur 15 Minuten für das Frühstück hat, mit einem 6-jährigen Kind und einer 60-jährigen Schwiegermutter zu Hause (\textbf{Hintergrund}). Bitte\textbf{liste eine 7-Tage-Frühstücks-Kombination ohne Wiederholung auf}(\textbf{Aufgabe}). Jedes Frühstück: (1) Zubereitungszeit max. 15 Minuten (2) nährstoffausgewogen (Eiweiß + Kohlenhydrate + Gemüse/Obst) (3) für Alt und Jung geeignet (4) Zutaten in normalen Supermärkten erhältlich (\textbf{Einschränkungen}).

\textbf{Warum wirksam}: Die Rolle legt die Fachlichkeit fest, der Hintergrund gibt die Gemeinsamkeit der drei Personengruppen, die Aufgabe definiert das Ergebnis, die 4 Einschränkungen sichern die Umsetzbarkeit.

\end{tcbinfo}

\begin{tcbinfo}{Szenario 2: Lernhilfe (Vorbereitung auf die Englisch-CET-6-Prüfung)}
\textbf{Prompt}:

Du bist ein\textbf{erfahrener Englisch-CET-6-Lehrer}(Rolle). Ich bin ein gewöhnlicher Bachelor-Student,\textbf{CET-4 mit 480 Punkten, CET-6-Ziel 500, noch 3 Monate}(Hintergrund). Bitte\textbf{erstelle für mich einen ausführbaren 12-Wochen-Vorbereitungsplan}(Aufgabe). Anforderungen: (1) Wöchentlich klare Zeitverteilung für Hören, Sprechen, Lesen und Schreiben (2) Schwerpunktrichtungen für jeden Monat angeben (3) 3 geeignete Lernmaterialien empfehlen (auf Chinesisch begründen) (4) In Tabellenform ausgeben (Einschränkung).

\textbf{Warum wirksam}: Die Rolle legt das Fachprofil fest, der Hintergrund gibt Punkte und Zeitanker, die Aufgabe verlangt einen „phasenweisen Plan“, die Einschränkungen fordern „Tabellenform“ und „Umsetzbarkeit“.

\end{tcbinfo}

\begin{tcbinfo}{Szenario 3: Arbeitseffizienz (Besprechungsprotokoll aufbereiten)}
\textbf{Prompt}:

Du bist ein\textbf{erfahrener Verwaltungsassistent}(Rolle). Ich habe gerade eine einstündige Abteilungsbesprechung beendet,\textbf{Thema war Q2-Marketing-Strategie}(Hintergrund).\textbf{Bitte fasse die folgenden Besprechungsnotizen in strukturierte Protokolle zusammen}(Aufgabe). Anforderungen: (1) Unterteilt in Beschlüsse / To-dos / Risiken und Einwände (2) To-dos mit Verantwortlichen und Fristen markieren (falls nicht erwähnt, “zu bestätigen”) (3) Gesamtzeichenanzahl unter 600 (4) In Markdown-Format ausgeben (Einschränkung).

\textbf{Warum wirksam}: A schreibt den Entwurf, B überarbeitet, C übersetzt – die Kombination ist noch wirkungsvoller.

\end{tcbinfo}

\textbf{Fortgeschrittene Anwendung: Lass die KI zu deinem “persönlichen Assistenten” werden}\par

System-Prompt

\begin{itemize}
  \item \textbf{Sieben. Grenzen und Herausforderungen der KI: Wachsame Urteilskraft bewahren}: Wenn du mit der Antwort der KI unzufrieden bist, kannst du sagen: "Das ist zu professionell, bitte erkläre es noch einmal in Worten, die ein Grundschüler versteht."
  \item \textbf{Chain-of-Thought (CoT)}: Wenn du mit der Antwort der KI unzufrieden bist, kannst du sagen: "Das ist zu professionell, bitte erkläre es noch einmal in Worten, die ein Grundschüler versteht."
  \item \textbf{KI ist ein Werkzeug, keine Autorität; bewahre eigenständiges Denken und Urteilsvermögen.}: A schreibt den ersten Entwurf, B feilt nach, C übersetzt – kombiniert wirkt es am besten.
  \item \textbf{Bei wichtigen Entscheidungen (medizinisch, juristisch, finanziell) konsultiere Fachleute anstatt dich auf KI zu verlassen.}: Lege in den KI-Einstellungen deine Identität, Präferenzen und Tabus fest, damit die Antworten besser zu dir passen.
\end{itemize}

\booksection{Sieben. Grenzen und Herausforderungen der KI: Wachsame Urteilskraft bewahren}

“Halluzinations”-Problem

\begin{itemize}
  \item \textbf{KI ist zwar mächtig, aber keineswegs "allmächtig". Ihre Grenzen nüchtern zu erkennen, ist Voraussetzung für ihren besseren Einsatz.}: KI kann ernsthaft Unsinn erzählen (Fachbegriff: Hallucination). Sie erfindet nicht existierende historische Ereignisse, falsche Zitate, gefälschte Statistiken. Wichtige Informationen müssen von Menschen überprüft werden.
  \item \textbf{Wissensschnitt}: Die Trainingsdaten der meisten KIs haben ein Stichtdatum (z. B. Anfang 2024); sie wissen nicht, was danach passiert ist.
  \item \textbf{Eingeschränkte Schlussfolgerungsfähigkeit}: Bei komplexen mehrstufigen Schlussfolgerungen, mathematischen Beweisen und logischen Fallen macht KI noch leicht Fehler.
  \item \textbf{Datenschutzrisiko}: Bei Cloud-KI-Diensten können deine Gespräche für das Modelltraining verwendet werden. Sensible Informationen sollten vorzugsweise mit On-Device-KI verarbeitet werden.
  \item \textbf{Voreingenommenheit und Werte}: KI kann Voreingenommenheit aus den Trainingsdaten aufweisen; eigenständiges Denken ist erforderlich.
\end{itemize}

\begin{tcbinfo}{KI-Nutzungstipps}
\begin{itemize}
  \item Von der KI erstellte Inhalte können ungenau sein ("Halluzination"); wichtige Informationen müssen überprüft werden.
  \item Gib keine persönlichen Daten in die KI ein (Personalausweisnummer, Passwörter, Bankdaten usw.).
  \item Nutze Prompt (Prompt) geschickt: Je konkreter und klarer die Beschreibung, desto höher die Qualität der KI-Antwort.
  \item Mehrstufige Dialogoptimierung
  \item Mehrere KIs zusammenarbeiten
\end{itemize}

\end{tcbinfo}

\begin{tcbpractice}{Herausforderung: Drei KI-Praxisaufgaben}
\begin{enumerate}
  \item \textbf{Herausforderung 1: Komplexe Anweisung}: Sage zum Sprachassistenten: "Erinnere mich morgen früh um 7 Uhr an den Regenschirm und verschiebe den heutigen 14-Uhr-Kalendereintrag auf 16 Uhr." Prüfe, ob er mehrere Schritte auf einmal ausführen kann.
  \item \textbf{Herausforderung 2: KI-Bildbearbeitung}: Fotografiere ein Touristenbild mit Passanten und nutze „magisches Entfernen“ oder den „KI-Radierer“, um die Passanten wegzureinigen. Beobachte, wie die KI den Hintergrund „ergänzt“.
  \item \textbf{Herausforderung 3: Prompt-Optimierung}: Frage die KI erst mit einem simplen Prompt wie „schreib einen Reiseführer“ und dann mit einem detaillierten Prompt nach dem CRISPE-Muster. Vergleiche die Qualität beider Antworten.
\end{enumerate}

\end{tcbpractice}

\begin{tcbinfo}{Fallbeispiel 2: Vollstaendiger Ablauf der Datenmigration beim Geraetewechsel}
\begin{itemize}
  \item \textbf{Szenario}: iPhone 13 Xiaomi 15 – markenuebergreifender Wechsel.
  \item \textbf{Schritte}: (1) Altgeraet: iCloud-/WeChat-eigenes Backup (2) Apple ID abmelden + „Suche" deaktivieren (3) Neues Geraet: „Xiaomi-Wechsel"-App scannen (4) Migrationsobjekte auswaehlen (5) WeChat separat mit „Chat-Verlaufs-Migration" (6) Schluesseldaten pruefen (Kontakte/Album/WeChat) (7) Altgeraet verschluesselt loeschen.
  \item \textbf{Erwartetes Ergebnis}: Kein Datenverlust; Altgeraet kann sicher weiterverkauft werden.
\end{itemize}

\end{tcbinfo}
